% 十、模型评价与推广
\section{模型评价与推广}
\subsection{模型的优点}
\begin{enumerate}
\item \textbf{四问一体的统一框架}：同一优化核、同一价值函数链（式\eqref{eq:chain}）、同一执行器贯穿四问，问题间只改变信息状态与交易权限；每个新增组件（正式预报、调整权限、价格场景）都有对应的消融或对照量化其价值，模型的递进是可验证的递进。
\item \textbf{理论结构支撑}：命题~\ref{prop:convex} 使价值函数以凸折线精确表示（无库存网格、无插值误差），命题~\ref{prop:readjust} 保证重优化方向正确，命题~\ref{prop:scale} 说明策略对价格水平结构稳健；三个命题均配独立数值验证。
\item \textbf{严格因果与无信息泄漏}：预测、场景、调参、执行各环节只用当时可得数据；全部参数由一月选定，评分期从未参与调参。
\item \textbf{结果可信度多重锚定}：LP 与连续 DP 双方法互证（750 组对照最大差 $3.64\times10^{-12}$ 元）；143 项自动化测试；全部 22 组全年实验由轨迹独立重算账单，库存递推、能量平衡与费用分解的复核误差均在 $10^{-12}$ 元量级；实验保存输入数据与代码指纹，可精确复现。
\item \textbf{按实际账单评价}：所有结论以真实交付数据结算的连续执行账单为准，并同时报告紧急购电与期末库存，代理目标与真实费用从不混淆。
\end{enumerate}

\subsection{模型的缺点与局限}
\begin{enumerate}
\item SAA 的场景内完美追索是乐观代理，本文未给出多阶段随机最优性的理论保证；平均未来费用函数 $\bar H$ 同样是机会成本近似而非精确随机价值函数。
\item 全部结论基于单一评分年（2025 年）数据，缺少跨年份外部验证；单因子消融收益不保证跨环境可加。
\item 经验残差场景假设“未来误差分布与近期历史相似”，对突变天气或负载结构变化的适应有滞后；场景等权而未做重要性加权。
\item 电价预测的日间水平模型较简单（回归 / AR(1)），极端价格跳变日（如 12 月 21 日）无法预见；所幸命题~\ref{prop:scale} 表明水平误差对决策形状影响有限。
\end{enumerate}

\subsection{改进方向与推广}
改进方向：(i) 对场景做基于天气类型的条件化选择或重要性加权，替代纯时间近邻；(ii) 价格水平模型引入外生温度/节假日特征；(iii) 用逼近动态规划在 $\bar H$ 上叠加误差修正项，压缩乐观代理与真实账单之间的口径差。推广方面，“信息状态—未来费用函数—因果执行”的框架不依赖本题细节：任何“提前承诺 + 库存缓冲 + 事后结算”的运营问题——电动车充电站的日前购电、冷库/蓄热系统的错峰用能、云计算的预留实例采购——都可以在同一骨架下实例化，只需替换预测层与结算式。凸折线价值函数的精确表示对一维库存状态普遍适用，计算量与时段数线性增长，可直接用于更细粒度（如 1 分钟）的执行控制。
